\documentclass[final]{beamer}

\usepackage[T1]{fontenc}
\usepackage[size=custom,width=120,height=90,scale=1.0]{beamerposter}
\usetheme{gemini}
\usecolortheme{labsix}

% Fonts
\usepackage{helvet}
\usepackage[helvet]{sfmath}
\renewcommand{\familydefault}{\sfdefault}

% Core packages
\usepackage{graphicx}
\usepackage{booktabs}
\usepackage[numbers]{natbib}
\usepackage{tikz}
\usetikzlibrary{positioning,arrows.meta,shapes.geometric,calc,fit,backgrounds}
\usepackage{xcolor}
\usepackage{array}
\usepackage{microtype}
\usepackage{ragged2e}
\usepackage{amsmath}

% Colors
\definecolor{myblue}{RGB}{64,115,158}
\definecolor{mydarkblue}{RGB}{39,60,117}
\definecolor{mylightblue}{RGB}{232,244,255}
\definecolor{mylightgray}{RGB}{245,246,250}
\definecolor{mygreen}{RGB}{40,130,90}
\definecolor{mylightgreen}{RGB}{220,248,235}

% Column geometry
\newlength{\sepwidth}
\newlength{\colwidth}
\setlength{\sepwidth}{0.020\paperwidth}
\setlength{\colwidth}{0.300\paperwidth}
\newcommand{\separatorcolumn}{\begin{column}{\sepwidth}\end{column}}

% Placeholder figure box
\newcommand{\figureplaceholder}[2][0.18\textheight]{%
  \begin{center}
    \fcolorbox{myblue}{white}{%
      \begin{minipage}[c][#1][c]{0.94\linewidth}
        \centering
        {\Large \textbf{Figure Placeholder}}\\[0.75em]
        {\normalsize Insert final figure here}
      \end{minipage}
    }\\[0.6em]
    {\small \textbf{#2}}
  \end{center}
}

\title{Age- and Arm-Stratified Co-expression Network Rewiring\\
in the Murine Kidney During and After Spaceflight}

\author{Ibrahim M.\ Shahid\textsuperscript{1,2,3}}

\institute{%
\textsuperscript{1}Middle College at UNCG \quad
\textsuperscript{2}Blue Marble Space Institute of Science \quad
\textsuperscript{3}NASA Human Research Program
}

\footercontent{%
\texttt{ibrahim.shahid@example.edu}
\hfill
ASGSR Annual Meeting 2026
\hfill
Dataset: NASA GeneLab OSD-771 (RRRM-2)
}

\begin{document}
\begin{frame}[t]
\begin{columns}[t]
\separatorcolumn

% ==========================================================
% COLUMN 1
% ==========================================================
\begin{column}{\colwidth}

\begin{block}{Background \& Objective}
Astronauts show elevated post-flight nephrolithiasis risk, and prior renal studies implicate distal nephron remodeling and transporter dysregulation during spaceflight. Standard differential expression (DE) analysis captures shifts in average transcript abundance, but it does not detect changes in how genes are coordinated with one another.

\medskip
This study asks whether spaceflight reorganizes \textbf{kidney co-expression network architecture} in a manner that depends on \textbf{age} and \textbf{flight arm} and whether these topological changes reveal biology that remains hidden to conventional DE analysis.

\medskip
We analyzed RRRM-2 / OSD-771, a full-factorial bulk kidney RNA-seq dataset, using a workflow that integrates:
\begin{itemize}
  \item nephron-segment deconvolution,
  \item confounder-aware residualization,
  \item shared-skeleton sample-specific network inference,
  \item node2vec embedding with Procrustes alignment,
  \item permutation-based rewiring significance testing, and
  \item pathway-level interpretation.
\end{itemize}
\end{block}

\begin{block}{Study Design}
\begin{center}
\renewcommand{\arraystretch}{1.35}
\begin{tabular}{>{\bfseries}l l}
Age & Young (16 wk) vs Old (34 wk) \\
Arm & ISS-T vs LAR \\
Environment Group & FLT, GC, VIV, BSL \\
Per cell & 5 mice \\
Total & 80 female C57BL/6NTac kidneys \\
\end{tabular}
\end{center}

\medskip
\textbf{Primary contrasts:}
ISS-T Young (FLT$-$GC), ISS-T Old (FLT$-$GC), LAR Young (FLT$-$GC), and LAR Old (FLT$-$GC).

\medskip
This design separates \textbf{acute in-flight remodeling} from \textbf{post-return recovery/persistence}, while also testing whether older kidneys respond differently from young kidneys.
\end{block}

\begin{block}{Analytical Workflow}
\textbf{Phase 0: Deconvolution} \\
MuSiC against a murine kidney single-cell reference to estimate nephron segment proportions.

\medskip
\textbf{Phase 1: Residualization} \\
DESeq2-VST expression with technical and compositional confounders removed while preserving biological structure of interest.

\medskip
\textbf{Phase 2: Shared topology} \\
A sparse global edge skeleton was learned after within-cell standardization to reduce Simpson's-paradox edge selection and improve cross-condition comparability.

\medskip
\textbf{Phase 3: Sample-specific networks} \\
LIONESS-style inference generated contrast-specific weighted networks on the same shared edge set.

\medskip
\textbf{Phase 4: Embedding and rewiring} \\
Node2vec embeddings were aligned by orthogonal Procrustes; per-gene rewiring was quantified as cosine-distance shift between matched contrasts.

\medskip
\textbf{Phase 5: Significance and interpretation} \\
Focused permutation testing prioritized highly rewired genes; enrichment analysis identified coherent biological programs.
\end{block}

\begin{block}{Recommended Figure: Pipeline Schematic}
\figureplaceholder[0.21\textheight]{Analytical pipeline integrating deconvolution, residualization, shared-skeleton network construction, LIONESS-based sample-specific weighting, node2vec embedding, Procrustes alignment, focused permutation testing, and pathway enrichment.}
\end{block}

\end{column}

\separatorcolumn

% ==========================================================
% COLUMN 2
% ==========================================================
\begin{column}{\colwidth}

\begin{block}{Genome-Wide Rewiring Is Broad and Orthogonal to DE}
Spaceflight produced broad remodeling of kidney co-expression structure across all four contrasts. Mean rewiring scores were similar across strata, indicating a \textbf{global transcriptomic reorganization} rather than a single isolated pathway effect.

\begin{center}
\renewcommand{\arraystretch}{1.2}
\begin{tabular}{lcc}
\toprule
\textbf{Contrast} & \textbf{Mean $\Delta$} & \textbf{Max $\Delta$} \\
\midrule
ISS-T Young & 0.284 & 0.515 \\
ISS-T Old   & 0.287 & 0.511 \\
LAR Young   & 0.271 & 0.513 \\
LAR Old     & 0.275 & 0.539 \\
\bottomrule
\end{tabular}
\end{center}

\medskip
Critically, rewiring magnitude showed \textbf{near-zero correlation with} $|\log_2 \mathrm{FC}|$, demonstrating that network topology captures a biological dimension that is largely independent of mean-expression change. In other words, genes can maintain similar abundance while being rewired into different regulatory neighborhoods.
\end{block}

\begin{alertblock}{Silent Shifters: Network Remodeling Hidden from DE}
We define \textbf{silent shifters} as genes with:
\begin{itemize}
  \item high rewiring (top decile),
  \item low absolute fold change, and
  \item no strong DE significance.
\end{itemize}

These genes are especially important because they suggest \textbf{regulatory restructuring without obvious expression-level change}. Across the four contrasts, the analysis identified dozens of silent shifters per condition, showing that a substantial fraction of the spaceflight response would be missed by standard transcriptomic screening alone.
\end{alertblock}

\begin{block}{Recommended Figure: Rewiring vs Differential Expression}
\figureplaceholder[0.21\textheight]{Rewiring versus $|\log_2\mathrm{FC}|$ across the four primary contrasts. Near-zero Spearman correlation supports the claim that network rewiring is largely orthogonal to classical differential expression.}
\end{block}

\begin{block}{Recommended Figure: Silent-Shifter Quadrant Plot}
\figureplaceholder[0.21\textheight]{Quadrant plot separating genes by statistical significance and rewiring magnitude. The high-rewiring / low-DE quadrant highlights silent shifter candidates that are biologically responsive despite minimal mean-expression change.}
\end{block}

\begin{block}{Focused Permutation Identifies Topological Hub Genes}
Focused permutation testing sharpened gene-level discovery by restricting multiple-testing burden to the most strongly rewired genes. The clearest signal occurred in \textbf{ISS-T Young}, where 11 genes achieved FDR $< 0.05$.

\medskip
\textbf{Top ISS-T Young hub hits included:}
\begin{itemize}
  \item \textbf{Slc6a20a}, \textbf{Slc22a28}: renal solute transport,
  \item \textbf{Trf}, \textbf{Gc}: secreted transport/homeostatic proteins,
  \item \textbf{Cry2}: circadian regulation,
  \item \textbf{Fabp7}, \textbf{Azgp1}, \textbf{Pfkp}: metabolic adaptation,
  \item \textbf{Foxq1}, \textbf{Adgrg4}: epithelial and signaling functions.
\end{itemize}

Older ISS-T animals showed a more restricted gene-level signature, led by \textbf{Prlr} and \textbf{Nr4a1}, consistent with altered hormonal and stress-response wiring.
\end{block}

\begin{block}{Recommended Figure: Gene-Rank / Focused Permutation Plot}
\figureplaceholder[0.20\textheight]{Focused-permutation rank plots for all four contrasts. ISS-T Young shows the strongest concentration of significant rewired hubs, while older and recovery-arm contrasts display sparser gene-level discovery.}
\end{block}

\end{column}

\separatorcolumn

% ==========================================================
% COLUMN 3
% ==========================================================
\begin{column}{\colwidth}

\begin{block}{Age- and Arm-Specific Pathway Remodeling}
Pathway enrichment revealed that the biological meaning of rewiring depends strongly on both \textbf{age} and \textbf{flight arm}.

\medskip
\textbf{ISS-T Young:}
ECM and complement-associated remodeling, consistent with acute structural and extracellular adaptation during flight.

\medskip
\textbf{ISS-T Old:}
EMT, steroid biosynthesis, ribosomal, and proteostasis-related signatures, consistent with a more stress-reactive and age-sensitive in-flight response.

\medskip
\textbf{LAR Young:}
PPAR-centered metabolic recovery programs after return to Earth.

\medskip
\textbf{LAR Old:}
PPAR, NOTCH, complement/coagulation, and translational pathways, suggesting a broader recovery state that still retains strong stress and remodeling components.

\medskip
Together, these results indicate that \textbf{young and old kidneys do not recover through the same transcriptomic topology}.
\end{block}

\begin{block}{Recommended Figure: Age-Dependent Rewiring}
\figureplaceholder[0.20\textheight]{Age-dependent rewiring comparison in ISS-T kidneys. Scatter and ranked-gene summaries highlight loci whose flight-associated network shifts diverge between young and old animals.}
\end{block}

\begin{block}{Recovery vs Persistence After Return}
Comparison of ISS-T and LAR contrasts separates rewiring that is \textbf{transient and flight-specific} from rewiring that \textbf{persists after return}. This helps distinguish acute spaceflight stress responses from longer-lasting network remodeling.

\medskip
A subset of highly rewired genes appears primarily in ISS-T, suggesting changes that resolve after re-entry. Another subset remains rewired in both arms, indicating more persistent systems-level remodeling.
\end{block}

\begin{block}{Recommended Figure: ISS-T vs LAR Recovery Plot}
\figureplaceholder[0.20\textheight]{ISS-T versus LAR rewiring comparison showing genes that remain above the diagonal only during spaceflight versus genes whose rewiring persists after return, thereby separating transient and persistent network effects.}
\end{block}

\begin{exampleblock}{Convergent Signal: Retinoid / Retinoic Acid Metabolism}
Retinoid metabolism emerged as the \textbf{most reproducible convergent signal} across multiple analytical layers.

\begin{itemize}
  \item At the pathway level, retinol metabolism was one of the strongest enrichment signals and peaked in the \textbf{LAR Young} recovery condition.
  \item At the gene level, several retinoid-associated genes showed consistently elevated rewiring across contrasts.
  \item Biologically, this pattern is more consistent with a \textbf{repair or recovery response} than with a purely acute in-flight lesion.
\end{itemize}

Because this signal appears most strongly after return, it provides a plausible mechanistic link between transcriptomic rewiring and renal recovery biology.
\end{exampleblock}

\begin{block}{Recommended Figure: Retinol / Retinoic Acid Panel}
\figureplaceholder[0.20\textheight]{Composite retinoid-metabolism figure showing pathway-level enrichment together with rewiring of representative retinol-pathway genes across contrasts, highlighting recovery-associated convergence in the LAR arm.}
\end{block}

\begin{alertblock}{Conclusions}
\begin{enumerate}
  \item Spaceflight induces \textbf{broad kidney co-expression network rewiring} across both ages and both flight arms.
  \item Rewiring is \textbf{largely orthogonal to differential expression}, so conventional DE alone misses a major component of the renal response.
  \item \textbf{Silent shifters} provide a concrete class of biologically responsive genes that change network context without large mean-expression shifts.
  \item The kidney response is \textbf{age-dependent}: younger and older animals remodel different pathway systems under flight and after return.
  \item \textbf{Retinoid metabolism} is a leading convergent signal and may reflect a recovery-associated remodeling program after spaceflight.
\end{enumerate}
\end{alertblock}

\begin{block}{References}
\footnotesize
[1] Kuijjer et al. \textit{iScience} 2019. \\
[2] Grover and Leskovec. \textit{KDD} 2016. \\
[3] Wang et al. \textit{Nat Commun} 2019. \\
[4] NASA GeneLab OSD-771 / RRRM-2. \\
[5] Relevant spaceflight kidney / nephrolithiasis references from your bibliography.
\end{block}

\end{column}

\separatorcolumn
\end{columns}
\end{frame}
\end{document}